\documentclass[10pt,a4paper]{amsart}
\usepackage[margin=19mm]{geometry}
\usepackage{amsmath,amssymb,mathtools}
\usepackage[scr=rsfs,cal=euler]{mathalfa}
\usepackage{xfrac}
\usepackage[unicode,hidelinks,bookmarks=false]{hyperref}
\newcommand{\ie}{i.e.\ }
\newcommand{\cf}{cf.\ }
\newcommand{\resp}{respectively\ }
\newcommand{\Subsection}{\subsection}
\providecommand{\nobreakdash}{\nobreak\hbox{-}}
\setlength{\emergencystretch}{2em}
\multlinegap0pt
\usepackage{bm}
\usepackage{bbm}
\usepackage{dsfont}
\usepackage{xspace}
\usepackage{cancel}
\usepackage{mathtools}
\usepackage{amsthm}
\makeatletter
\@ifundefined{theorem}{\newtheorem{theorem}{Theorem}}{}
\@ifundefined{lemma}{\newtheorem{lemma}[theorem]{Lemma}}{}
\@ifundefined{theoremalph}{\newtheorem{theoremalph}{Theorem}}{}
\makeatother
\newcommand{\R}{\mathbb{R}}
\title[Existence of extremal functions and Wulff symmetry for aniso]{Existence of extremal functions and Wulff symmetry for anisotropic Trudinger-Moser inequalities}
\author{Kaiwen Guo, Yanjun Liu}
\date{Source: arXiv:2511.04342v3 (CC BY 4.0)}
\begin{document}
\maketitle

\noindent\emph{Source and licence.} Existence of extremal functions and Wulff symmetry for anisotropic Trudinger-Moser inequalities. Version arXiv:2511.04342v3 (CC BY 4.0), distributed under the Creative Commons Attribution 4.0 International licence, as recorded in the arXiv licence metadata for this exact version and shown on the abstract page https://arxiv.org/abs/2511.04342v3. Full rights notice: licenses/arxiv-2511.04342-CC-BY-4.0.txt.\par\smallskip

\noindent\textit{Editorial scope.} Counted unit: the Lemma whose source label is \texttt{le2.2} in the article (source0.tex lines 289--294, source label \texttt{le2.2}), delivered together with the article's own complete original proof of it (source0.tex lines 392--536, one \texttt{proof} environment of 145 lines). No mathematics is written, repaired or re-derived: statement and proof are the article's own text line for line. The proof is detached from the statement in the source: the article states the result at lines 289--294 and prints its proof at lines 392--536, and the proof environment's own header names the statement, so the association is the article's own. Required context retained uncounted: tha (lines 172--178); le2.3 (lines 295--301); le2.4 (lines 305--314); le2.5 (lines 331--337). Every definition, display and label of those lines is retained unchanged. Omitted from this package and not redistributed: 1--171, 179--288, 302--304, 315--330, 338--391, 537--821. Nothing inside a retained excerpt is omitted or altered: every retained body is delivered exactly as the article's lines are written, so no omission receipt of any kind accompanies this package. Citations: the retained excerpts carry no citation command at all, so no bibliography entry of the article is transcribed and no reference is redistributed. Reference adaptation: none was needed and none was made; every reference inside the delivered unit resolves inside it, so no unresolved reference or citation is produced. Printed slips kept literal: no printed word, symbol, hypothesis or number of the counted statement or of its proof was altered. Layout and class: the article's own class is \texttt{article} with packages including CJK, amsmath, amsfonts, amssymb, amsthm, amscd, amsbsy, mathrsfs, fancyhdr, geometry, graphicx, enumerate, indentfirst, color; the wrapper is the lane's pinned \texttt{amsart} editorial wrapper at 10pt on A4, which restates the theorem-like environments the retained text needs and adds the article's own macro definitions that the retained text needs (\texttt{\textbackslash R}), with extra wrapper packages (bm, bbm, dsfont, xspace, cancel, mathtools). The delivered numbering is the unit's own. Assets: no retained span contains a figure environment, a graphics-inclusion command, a PDF-page inclusion, a table or a TikZ picture; no image file is copied into this package, the package declares no asset dependency and no image file is redistributed with it. Licence: version arXiv:2511.04342v3 of this article is distributed under the Creative Commons Attribution 4.0 International licence, as recorded in the arXiv licence metadata for this exact version and shown on the abstract page https://arxiv.org/abs/2511.04342v3. A full rights notice is delivered at licenses/arxiv-2511.04342-CC-BY-4.0.txt. Characters: every retained excerpt is pure ASCII, so the delivered engine, XeLaTeX with its default Latin Modern text font, reports no missing character.
\begin{theoremalph}(Sharp subcritical anisotropic Trudinger-Moser inequality).\label{tha}
Let $p>q(1-\frac{\beta}{N})$ if $\beta\neq0$ and $p\geq q$ if $\beta=0$. Then
\begin{align}
    \mathrm{ATMSC}(N,p,q,\lambda,\beta)&:=\sup\limits_{\|F(\nabla u)\|_{N}\leq 1}\frac{1}{\|u\|_{q}^{q(1-\frac{\beta}{N})}}\int\limits_{\R^{N}}\frac{\exp(\lambda(1-\frac{\beta}{N})\lvert u\rvert^{\frac{N}{N-1}})\lvert u\rvert^{p}}{F^{o}(x)^{\beta}}\;\mathrm{d}x<+\infty,\notag\\
    \mathrm{ATMSC}(N,q,\lambda,\beta)&:=\sup\limits_{\|F(\nabla u)\|_{N}\leq 1}\frac{1}{\|u\|_{q}^{q(1-\frac{\beta}{N})}}\int\limits_{\R^{N}}\frac{\Phi_{N,q,\beta}(\lambda(1-\frac{\beta}{N})\lvert u\rvert^{\frac{N}{N-1}})}{F^{o}(x)^{\beta}}\;\mathrm{d}x<+\infty,\notag
\end{align}
\end{theoremalph}

\begin{lemma}\label{le2.3}(Hardy-Littlewood Inequality)
Let $f$, $g$ be nonnegative functions on $\R^{N}$, vanishing at infinity. Then
$$
\int\limits_{\R^{N}}f(x)g(x)\;\mathrm{d}x\leq \int\limits_{\R^{N}}f^{\star}(x)g^{\star}(x)\;\mathrm{d}x.
$$
Moreover, if $f$ is strictly decreasing and Wulff symmetric, then there is equality if and only if $g=g^{\star}$.
\end{lemma}

\begin{lemma}\label{le2.4}
Let $\Omega\subset\R^{N}$, $\lvert\Omega\rvert<+\infty$. Assume that
$$
f_{n}\rightarrow f\;a.e.\;in\;\Omega
$$
and there exists $q>1$ such that $f_{n}$ is uniformly bounded in $L^{q}(\Omega)$, $f\in L^{q}(\Omega)$. Then
$$
f_{n}\rightarrow f\;in\;L^{1}(\Omega).
$$
\end{lemma}

\begin{lemma}\label{le2.5}
Let $p>q(1-\frac{\beta}{N})$ if $\beta\neq0$ and $p\geq q$ if $\beta=0$. Then
\begin{align}
    \mathrm{ATMSC}(N,p,q,\lambda,\beta)&=\sup\limits_{\|F(\nabla u)\|_{N}=\|u\|_{q}=1}\int\limits_{\R^{N}}\frac{\exp(\lambda(1-\frac{\beta}{N})\lvert u\rvert^{\frac{N}{N-1}})\lvert u\rvert^{p}}{F^{o}(x)^{\beta}}\;\mathrm{d}x,\notag\\
    \mathrm{ATMSC}(N,q,\lambda,\beta)&=\sup\limits_{\|F(\nabla u)\|_{N}=\|u\|_{q}=1}\int\limits_{\R^{N}}\frac{\Phi_{N,q,\beta}(\lambda(1-\frac{\beta}{N})\lvert u\rvert^{\frac{N}{N-1}})}{F^{o}(x)^{\beta}}\;\mathrm{d}x.\notag
\end{align}
\end{lemma}

\begin{lemma}\label{le2.2}(P$\acute{o}$lya-Szeg$\ddot{o}$ Inequality)
Let $u\in D^{N,q}(\R^{N})$. Then $u^{\star}\in D^{N,q}(\R^{N})$ and
\begin{align}
    &\int\limits_{\R^{N}}F(\nabla u^{\star})^{N}\;\mathrm{d}x   \leq\int\limits_{\R^{N}}F(\nabla u)^{N}\;\mathrm{d}x.\notag
\end{align}
\end{lemma}

\begin{proof}
Let $\{\omega_{n}\}$ be a maximizing sequence of
$$
\mathrm{ATMSC}(N,p,q,\lambda,\beta)=\sup\limits_{\|F(\nabla u)\|_{N}\leq 1}\frac{1}{\|u\|_{q}^{q(1-\frac{\beta}{N})}}\int\limits_{\R^{N}}\frac{\exp(\lambda(1-\frac{\beta}{N})\lvert u\rvert^{\frac{N}{N-1}})\lvert u\rvert^{p}}{F^{o}(x)^{\beta}}\;\mathrm{d}x.
$$
Then by Lemma \ref{le2.2}, Lemma \ref{le2.3} and Lemma \ref{le2.5}, we may assume that $\omega_{n}$ is nonnegative, Wulff symmetric and
$$
\int\limits_{\R^{N}}\frac{\exp(\lambda(1-\frac{\beta}{N})\lvert \omega_{n}\rvert^{\frac{N}{N-1}})\lvert \omega_{n}\rvert^{p}}{F^{o}(x)^{\beta}}\;\mathrm{d}x\rightarrow\mathrm{ATMSC}(N,p,q,\lambda,\beta)
$$
with
$$
\|F(\nabla\omega_{n})\|_{N}=\|\omega_{n}\|_{q}=1.
$$
For $q>1$, up to a subsequence, there exists some $\omega\in D^{N,q}(\R^{N})$ such that
\begin{align}
    &\omega_{n}\rightharpoonup \omega\;\;\mathrm{weakly}\;\mathrm{in}\;D^{N,q}(\R^{N})\textup{;}\notag\\
    &\omega_{n}\rightarrow \omega\;\;\mathrm{strongly}\;\mathrm{in}\;L^{N}(\R^{N})\;\mathrm{and}\;L_{\mathrm{loc}}^{i}(\R^{N})\;\mathrm{for}\;i\in(0,+\infty)\textup{;}\notag\\
    &\omega_{n}\rightarrow \omega\;\;\mathrm{a.e.}\;\mathrm{in}\;\R^{N}.\notag
\end{align}
Obviously, $\omega$ is nonnegative, Wulff symmetric and
$$
\|F(\nabla\omega)\|_{N}\leq1,\;\|\omega\|_{q}\leq1.
$$

{\bfseries Case 1:} $\beta=0$, $p>q$.

For any $x\not=0$,
\begin{align}
    \|\omega_{n}\|_{q}^{q}\geq\int_{0}^{F^{o}(x)}\omega_{n}(r)^{q}\int_{\partial W_{r}}\frac{1}{\lvert\nabla F^{o}(x)\rvert}\;\mathrm{d}\sigma\mathrm{d}r\geq\kappa_{N}F^{o}(x)^{N}\omega_{n}(F^{o}(x))^{q}.\notag
\end{align}
Thus, for any $\epsilon>0$, there exists $R>0$ sufficiently large such that $\lvert\omega_{n}(x)\rvert\leq\epsilon$ when $F^{o}(x)\geq R$. Then by Theorem \ref{tha},
\begin{align}
    \int\limits_{\R^{N}\backslash W_{R}}\exp(\lambda\lvert \omega_{n}\rvert^{\frac{N}{N-1}})\lvert \omega_{n}\rvert^{p}\;\mathrm{d}x&\leq\epsilon^{p-q}\int\limits_{\R^{N}\backslash W_{R}}\exp(\lambda\lvert \omega_{n}\rvert^{\frac{N}{N-1}})\lvert \omega_{n}\rvert^{q}\;\mathrm{d}x\notag\\
    &\leq C(N,q,\lambda)\epsilon^{p-q}.\notag
\end{align}
Next, notice that
$$
\exp(\lambda\lvert \omega_{n}\rvert^{\frac{N}{N-1}})\lvert \omega_{n}\rvert^{p}\rightarrow\exp(\lambda\lvert \omega\rvert^{\frac{N}{N-1}})\lvert \omega\rvert^{p}\;\mathrm{a.e.}\;\mathrm{in}\;W_{R}
$$
and by Theorem \ref{tha}, there exists $\theta>0$ sufficiently small such that $\exp(\lambda\lvert \omega_{n}\rvert^{\frac{N}{N-1}})\lvert \omega_{n}\rvert^{p}$ is uniformly bounded in $L^{1+\theta}(W_{R})$. Thus, by Lemma \ref{le2.4},
$$
\int\limits_{W_{R}}\exp(\lambda\lvert \omega_{n}\rvert^{\frac{N}{N-1}})\lvert \omega_{n}\rvert^{p}\;\mathrm{d}x\rightarrow\int\limits_{W_{R}}\exp(\lambda\lvert \omega\rvert^{\frac{N}{N-1}})\lvert \omega\rvert^{p}\;\mathrm{d}x.
$$
Hence,
\begin{align}
    \mathrm{ATMSC}(N,p,q,\lambda,\beta)\leftarrow&\int\limits_{\R^{N}}\exp(\lambda\lvert \omega_{n}\rvert^{\frac{N}{N-1}})\lvert \omega_{n}\rvert^{p}\;\mathrm{d}x\notag\\
    \leq&\int\limits_{\R^{N}}\exp(\lambda\lvert \omega\rvert^{\frac{N}{N-1}})\lvert \omega\rvert^{p}\;\mathrm{d}x+C(N,q,\lambda)\epsilon^{p-q}.\notag
\end{align}
Since $\epsilon>0$ is arbitrary, we obtain
$$
\mathrm{ATMSC}(N,p,q,\lambda,\beta)\leq\int\limits_{\R^{N}}\exp(\lambda\lvert \omega\rvert^{\frac{N}{N-1}})\lvert \omega\rvert^{p}\;\mathrm{d}x.
$$
Obviously, $\omega\not=0$ and
\begin{align}
    \mathrm{ATMSC}(N,p,q,\lambda,\beta)&\leq\int\limits_{\R^{N}}\exp(\lambda\lvert \omega\rvert^{\frac{N}{N-1}})\lvert \omega\rvert^{p}\;\mathrm{d}x\notag\\
    &\leq\frac{1}{\|\omega\|_{q}^{q}}\int\limits_{\R^{N}}\exp(\lambda\lvert \omega\rvert^{\frac{N}{N-1}})\lvert \omega\rvert^{p}\;\mathrm{d}x.\notag
\end{align}
Therefore, $\|\omega\|_{q}=1$ and
$$
\mathrm{ATMSC}(N,p,q,\lambda,\beta)=\int\limits_{\R^{N}}\exp(\lambda\lvert \omega\rvert^{\frac{N}{N-1}})\lvert \omega\rvert^{p}\;\mathrm{d}x.
$$
Moreover, $\omega$ is a maximizer of $\mathrm{ATMSC}(N,p,q,\lambda,\beta)$ with $\|F(\nabla\omega)\|_{N}\leq1$.

{\bfseries Case 2:} $\beta=0$, $p=q$.

Similarly as in the first case, we could show that for any $\epsilon>0$, there exists $R>0$ sufficiently large such that $\lvert\omega_{n}(x)\rvert\leq\epsilon$ when $F^{o}(x)\geq R$ and
$$
\int\limits_{W_{R}}\left[\exp(\lambda\lvert \omega_{n}\rvert^{\frac{N}{N-1}})-1\right]\lvert \omega_{n}\rvert^{q}\;\mathrm{d}x\rightarrow\int\limits_{W_{R}}\left[\exp(\lambda\lvert \omega\rvert^{\frac{N}{N-1}})-1\right]\lvert \omega\rvert^{q}\;\mathrm{d}x.
$$
Meanwhile, by Theorem \ref{tha} and $e^{t}\leq1+te^{t}$ for $t\geq0$,
\begin{align}
    \int\limits_{\R^{N}\backslash W_{R}}\left[\exp(\lambda\lvert \omega_{n}\rvert^{\frac{N}{N-1}})-1\right]\lvert \omega_{n}\rvert^{q}\;\mathrm{d}x&\leq\lambda\int\limits_{\R^{N}\backslash W_{R}}\exp(\lambda\lvert \omega_{n}\rvert^{\frac{N}{N-1}})\lvert \omega_{n}\rvert^{q+\frac{N}{N-1}}\;\mathrm{d}x\notag\\
    &\leq C(N,q,\lambda)\epsilon^{\frac{N}{N-1}}.\notag
\end{align}
Thus,
\begin{align}
    \mathrm{ATMSC}(N,p,q,\lambda,\beta)\leftarrow&\int\limits_{\R^{N}}\exp(\lambda\lvert \omega_{n}\rvert^{\frac{N}{N-1}})\lvert \omega_{n}\rvert^{q}\;\mathrm{d}x\notag\\
    =&\int\limits_{\R^{N}}\left[\exp(\lambda\lvert \omega_{n}\rvert^{\frac{N}{N-1}})-1\right]\lvert \omega_{n}\rvert^{q}\;\mathrm{d}x+\|\omega_{n}\|_{q}^{q}\notag\\
    \leq&\int\limits_{\R^{N}}\left[\exp(\lambda\lvert \omega\rvert^{\frac{N}{N-1}})-1\right]\lvert \omega\rvert^{q}\;\mathrm{d}x+C(N,q,\lambda)\epsilon^{\frac{N}{N-1}}+1.\notag
\end{align}
Since $\epsilon>0$ is arbitrary, we obtain
$$
\mathrm{ATMSC}(N,p,q,\lambda,\beta)\leq\int\limits_{\R^{N}}\left[\exp(\lambda\lvert \omega\rvert^{\frac{N}{N-1}})-1\right]\lvert \omega\rvert^{q}\;\mathrm{d}x+1.
$$
Obviously, $\omega\not=0$ and
\begin{align}
    \mathrm{ATMSC}(N,p,q,\lambda,\beta)&\leq\frac{1}{\|\omega\|_{q}^{q}}\int\limits_{\R^{N}}\left[\exp(\lambda\lvert \omega\rvert^{\frac{N}{N-1}})-1\right]\lvert \omega\rvert^{q}\;\mathrm{d}x+1\notag\\
    &=\frac{1}{\|\omega\|_{q}^{q}}\int\limits_{\R^{N}}\exp(\lambda\lvert \omega\rvert^{\frac{N}{N-1}})\lvert \omega\rvert^{q}\;\mathrm{d}x.\notag
\end{align}
Therefore, $\|\omega\|_{q}=1$ and
$$
\mathrm{ATMSC}(N,p,q,\lambda,\beta)=\int\limits_{\R^{N}}\exp(\lambda\lvert \omega\rvert^{\frac{N}{N-1}})\lvert \omega\rvert^{q}\;\mathrm{d}x.
$$
Moreover, $\omega$ is a maximizer of $\mathrm{ATMSC}(N,p,q,\lambda,\beta)$ with $\|F(\nabla\omega)\|_{N}\leq1$.

{\bfseries Case 3:} $0<\beta<N$.

Similarly as in the first case, we could show that for any $\epsilon>0$, there exists $R>0$ sufficiently large such that $\lvert\omega_{n}(x)\rvert\leq\epsilon$ when $F^{o}(x)\geq R$ and
$$
\int\limits_{W_{R}}\frac{\exp(\lambda(1-\frac{\beta}{N})\lvert \omega_{n}\rvert^{\frac{N}{N-1}})\lvert \omega_{n}\rvert^{p}}{F^{o}(x)^{\beta}}\;\mathrm{d}x\rightarrow\int\limits_{W_{R}}\frac{\exp(\lambda(1-\frac{\beta}{N})\lvert \omega\rvert^{\frac{N}{N-1}})\lvert \omega\rvert^{p}}{F^{o}(x)^{\beta}}\;\mathrm{d}x.
$$
Meanwhile, for $\delta=\delta(N,p,q,\beta)>0$ sufficiently small, by Theorem \ref{tha},
\begin{align}
    \int\limits_{\R^{N}\backslash W_{R}}\frac{\exp(\lambda(1-\frac{\beta}{N})\lvert \omega_{n}\rvert^{\frac{N}{N-1}})\lvert \omega_{n}\rvert^{p}}{F^{o}(x)^{\beta}}\;\mathrm{d}x&\leq\frac{1}{R^{\delta}}\int\limits_{\R^{N}\backslash W_{R}}\frac{\exp(\lambda(1-\frac{\beta}{N})\lvert \omega_{n}\rvert^{\frac{N}{N-1}})\lvert \omega_{n}\rvert^{p}}{F^{o}(x)^{\beta-\delta}}\;\mathrm{d}x\notag\\
    &\leq C(N,p,q,\lambda,\beta)\frac{1}{R^{\delta}}.\notag
\end{align}
Thus,
\begin{align}
    \mathrm{ATMSC}(N,p,q,\lambda,\beta)\leftarrow&\int\limits_{\R^{N}}\frac{\exp(\lambda(1-\frac{\beta}{N})\lvert \omega_{n}\rvert^{\frac{N}{N-1}})\lvert \omega_{n}\rvert^{p}}{F^{o}(x)^{\beta}}\;\mathrm{d}x\notag\\
    \leq&\int\limits_{\R^{N}}\frac{\exp(\lambda(1-\frac{\beta}{N})\lvert \omega\rvert^{\frac{N}{N-1}})\lvert \omega\rvert^{p}}{F^{o}(x)^{\beta}}\;\mathrm{d}x+C(N,p,q,\lambda,\beta)\frac{1}{R^{\delta}}.\notag
\end{align}
Since $R>0$ is sufficiently large, we obtain
$$
\mathrm{ATMSC}(N,p,q,\lambda,\beta)\leq\int\limits_{\R^{N}}\frac{\exp(\lambda(1-\frac{\beta}{N})\lvert \omega\rvert^{\frac{N}{N-1}})\lvert \omega\rvert^{p}}{F^{o}(x)^{\beta}}\;\mathrm{d}x.
$$
Obviously, $\omega\not=0$ and
\begin{align}
    \mathrm{ATMSC}(N,p,q,\lambda,\beta)&\leq\int\limits_{\R^{N}}\frac{\exp(\lambda(1-\frac{\beta}{N})\lvert \omega\rvert^{\frac{N}{N-1}})\lvert \omega\rvert^{p}}{F^{o}(x)^{\beta}}\;\mathrm{d}x\notag\\
    &\leq\frac{1}{\|\omega\|_{q}^{q(1-\frac{\beta}{N})}}\int\limits_{\R^{N}}\frac{\exp(\lambda(1-\frac{\beta}{N})\lvert \omega\rvert^{\frac{N}{N-1}})\lvert \omega\rvert^{p}}{F^{o}(x)^{\beta}}\;\mathrm{d}x.\notag
\end{align}
Therefore, $\|\omega\|_{q}=1$ and
$$
\mathrm{ATMSC}(N,p,q,\lambda,\beta)=\int\limits_{\R^{N}}\frac{\exp(\lambda(1-\frac{\beta}{N})\lvert \omega\rvert^{\frac{N}{N-1}})\lvert \omega\rvert^{p}}{F^{o}(x)^{\beta}}\;\mathrm{d}x.
$$
Moreover, $\omega$ is a maximizer of $\mathrm{ATMSC}(N,p,q,\lambda,\beta)$ with $\|F(\nabla\omega)\|_{N}\leq1$.

Finally, we prove that $\|F(\nabla\omega)\|_{N}=1$ when $\omega$ is a maximizer of $\mathrm{ATMSC}(N,p,q,\lambda,\beta)$. We assume $\|F(\nabla\omega)\|_{N}=\gamma^{-1}$ with $\gamma\geq1$, where we can easily exclude the case $\|F(\nabla u)\|_{N}=0$. Let
$$
v(x)=\gamma\omega(\gamma^{\frac{q}{N}}x).
$$
Obviously,
$$
\|F(\nabla v)\|_{N}=1,\;\|v\|_{q}=\|\omega\|_{q}
$$
and
\begin{align}
    &\frac{1}{\|v\|_{q}^{q(1-\frac{\beta}{N})}}\int\limits_{\R^{N}}\frac{\exp(\lambda(1-\frac{\beta}{N})\lvert v\rvert^{\frac{N}{N-1}})\lvert v\rvert^{p}}{F^{o}(x)^{\beta}}\;\mathrm{d}x\notag\\
    =&\gamma^{p-q(1-\frac{\beta}{N})}\frac{1}{\|\omega\|_{q}^{q(1-\frac{\beta}{N})}}\int\limits_{\R^{N}}\frac{\exp(\gamma^{\frac{N}{N-1}}\lambda(1-\frac{\beta}{N})\lvert \omega\rvert^{\frac{N}{N-1}})\lvert \omega\rvert^{p}}{F^{o}(x)^{\beta}}\;\mathrm{d}x\notag\\
    \geq&\frac{1}{\|\omega\|_{q}^{q(1-\frac{\beta}{N})}}\int\limits_{\R^{N}}\frac{\exp(\lambda(1-\frac{\beta}{N})\lvert \omega\rvert^{\frac{N}{N-1}})\lvert \omega\rvert^{p}}{F^{o}(x)^{\beta}}\;\mathrm{d}x\notag\\
    =&\mathrm{ATMSC}(N,p,q,\lambda,\beta)\notag.
\end{align}
Therefore, $\gamma=1$.

Maximizers for $\mathrm{ATMSC}(N,q,\lambda,\beta)$ can be proved similarly.
\end{proof}

\end{document}
